\chapter{Introduction}
\label{chap:introduction}

As a principal data source for perception in autonomous driving, \ac{LiDAR}
provides direct and precise range measurements of the three-dimensional
environment~\cite{geiger2013kitti}. The measurement process is, however, inherently non-uniform.
Because a \ac{LiDAR} sensor emits rotating laser beams at a fixed angular
resolution, two adjacent beams are always separated by a constant angle:
close to the sensor the beams lie near one another, whereas their separation
grows steadily with distance. An object of a given size therefore receives
progressively fewer laser returns the farther away it is located. As a
consequence, \ac{LiDAR} point clouds are dense in the near range and sparse
in the far range, and they additionally suffer from occlusion and from an
insufficient number of points on small objects. The density of the points in
turn affects both the classification of an object and the localisation of
its three-dimensional bounding box. In the KITTI detection
benchmark~\cite{geiger2012kitti,geiger2013kitti}, objects that are more
distant or occluded and thus carry fewer points are harder to detect.

Precisely because the detection difficulty is caused by a shortage of
points, an idea suggests itself: to make the points denser, that is, to
generate a denser and more uniform point cloud from the sparse input, in the
hope of recovering the object surface and its local structure. In other
words, before the data enter the classification, segmentation or detection
module, an independent module could generate additional points that lie on
the observed surfaces and thereby recover part of the lost performance ---
without replacing the sensor and without retraining the downstream detector.
Point cloud upsampling provides exactly such methods. Covering both
object-level and scene-level problems, this thesis compares different
upsampling methods, comprising classical geometric-optimisation approaches,
represented by edge-aware resampling~\cite{huang2013ear}, and
deep-learning-based
approaches~\cite{yu2018punet,qian2021pugcn,kim2023puedgeformer,%
zhang2025pdans}. Whether a method that performs better on geometric metrics
also performs better in downstream classification and detection, however,
remains open. This leads to the fundamental question addressed in this
thesis.

To examine this relationship, it is necessary to distinguish geometric
reconstruction quality from downstream task performance. Geometric
evaluation in point cloud upsampling uses distances to reference points or
surfaces and measures of sampling uniformity~\cite{yu2018punet,qian2021pugcn}.
These include Chamfer distance (CD), Hausdorff distance (HD),
point-to-surface distance and uniformity measures. They can reflect changes
in surface boundaries and local point distribution, so these properties
are not exclusive to detector evaluation. However, such geometric scores
do not directly evaluate a downstream network's predicted classes or
bounding boxes. PointRCNN uses foreground classification and proposal
refinement, whereas CenterPoint predicts object centres and box
attributes~\cite{shi2019pointrcnn,yin2021centerpoint}. Whether an
improvement in geometric reconstruction also improves these predictions
therefore needs to be evaluated on the downstream task.

The automotive setting poses a further difficulty when applying networks
trained on sampled object surfaces to \ac{LiDAR} scenes. PU-Net, for
example, learns from point patches sampled from 3D
models~\cite{yu2018punet}. This differs from the partial observations
obtained by a vehicle-mounted \ac{LiDAR}
sensor~\cite{geiger2012kitti,geiger2013kitti}. In this thesis, the
integration of the selected patch-based networks includes local patch
extraction and normalisation. In the later PU-GCN pipeline, predicted
patches are rescaled and translated back to \ac{LiDAR} coordinates before
merging and detector-input construction
(Section~\ref{sec:fc-4-5-2}). Normalisation is therefore reversed before
the downstream task; it does not itself imply a permanent change in scene
scale. The integration experiments examine coordinate-transformation
errors and the effects of patch construction and merging
(Section~\ref{sec:fc-4-12}). This allows implementation effects to be
examined separately from the behaviour attributed to the upsampling method.

This thesis therefore investigates upsampling under controlled output
point counts and explicitly specified downstream training conditions.
The central question is whether upsampling recovers useful surface
evidence and improves downstream classification or detection. A related
question is how changes in geometric quality correspond to changes in
these task outcomes. The study compares CD and HD with downstream
performance and considers the recorded point-to-face (P2F) distances as a
separate surface diagnostic. The different normalisation used for this
diagnostic is specified in Section~\ref{sec:geometry-audit}.

To address these questions, this thesis makes the following contributions.

\begin{enumerate}
    \item \textbf{A strict, cardinality-controlled evaluation protocol.}
    Within each input line, upsampling methods produce the same target
    number of output points for a given input. Original and sparse
    baselines retain their native point counts. Direct outputs are
    distinguished from observed-first KITTI inputs, which concatenate
    all measured input rows with selected predicted rows. Preservation
    applies to the constructed point-cloud file in that condition; direct
    outputs do not impose this requirement. Later detector sampling or
    voxelisation can still limit which points enter the network
    (Sections~\ref{sec:fc-4-6} and~\ref{sec:fc-4-8}).

    \item \textbf{A two-line experimental design.} Line~A tests whether
    densifying the original input improves performance over that original
    input. Line~B tests recovery after deliberate fourfold downsampling,
    using the sparse input as the baseline and the original input as a
    recovery reference. Results are reported separately because the lines
    start from different observations and address different performance
    questions.

    \item \textbf{A cross-task evaluation framework.} The same protocol is
    applied both to object-level shape classification on
    ModelNet40~\cite{wu2015modelnet} with PointNet++~\cite{qi2017pointnetpp}
    and to scene-level 3D detection on KITTI --- with the point-based
    detector PointRCNN~\cite{shi2019pointrcnn} and the voxel-based detector
    CenterPoint~\cite{yin2021centerpoint}. The two detectors consume the
    input in structurally different ways, which makes it possible to
    distinguish the effect of the generated geometry itself from the effect
    of a particular detector architecture; the object-level branch, in turn,
    provides a controlled surface-sampling setting in which to examine
    whether the same behaviour also occurs outside real sensor scenes.

    \item \textbf{Task-level evaluation.} Geometric metrics are used
    together with the classification accuracy of PointNet++ and the
    detection \ac{AP} of PointRCNN and CenterPoint.

    \item \textbf{An explicit treatment of the input budget.} This thesis
    formalises the observation that a detector accepts only a bounded
    amount of input, so that adding points beyond that budget forces a selection. The protocol includes configurations that separate
    this effect from the geometric quality of the generated points.

    \item \textbf{The integration process itself as an object of study.}
    The spatial locality of the patches, unique support, coverage and
    normalisation are explicitly defined, measured and audited rather than
    taken for granted, so that conclusions about a method do not
    inadvertently become conclusions about its wrapper code.
\end{enumerate}

To describe this evaluation design consistently, the experiments
distinguish between \emph{task branches} and \emph{input lines}.

Two task branches are studied:
\begin{itemize}
    \item classification with PointNet++ on ModelNet40, and
    \item detection with PointRCNN and CenterPoint on KITTI.
\end{itemize}

Each task branch contains two input lines: Line~A densifies the original
observation, whereas Line~B first degrades the observation
by a deterministic fourfold reduction and then attempts to recover it. Both
lines are specified formally in Section~\ref{sec:fc-3-3}.
Throughout this thesis the terminology is therefore used uniformly as
follows: \emph{task branch} refers to ModelNet40 or KITTI, and \emph{input
line} refers to Line~A or Line~B.

The primary detection comparison keeps the downstream detectors frozen.
A later, separate experiment adapts PointRCNN and CenterPoint to PU-GCN
inputs while retaining a fixed PU-GCN upsampler. This is detector adaptation,
not retraining of the upsampling network. The object-classification branch
uses separately trained PointNet++ SSG classifiers under the common training
recipe described in Chapter~\ref{chap:setup}; it is distinct from this later
detector-adaptation experiment.

The following chapters describe the foundations, experimental design and
results of this evaluation. Chapter~\ref{chap:fundamentals} introduces
point clouds, \ac{LiDAR} perception and the networks used for classification
and detection. It explains the upsampling methods and their mathematical
formulation. It also defines the geometric and task metrics used in the
evaluation.
Chapter~\ref{chap:concept} develops the concept: the research hypothesis,
the evaluation pipeline, the two comparison lines, the fairness protocol,
and the challenges that had to be solved. Chapter~\ref{chap:setup} documents
the experimental setup at the level of detail required for reproduction:
data preparation, downsampling, the integration of each upsampling method,
the format conversions, the detector and classifier configurations, the
computation of the geometric metrics, and the verification procedures.
Chapter~\ref{chap:results} presents and discusses the results.
Chapter~\ref{chap:conclusion} summarises the findings and outlines the
direction in which an upsampling method would have to be developed in order
to be genuinely useful for downstream 3D perception.
